\documentclass[reqno,11pt]{amsart}
\usepackage{math327-lecture}
\usepackage{needspace}
\usepackage{tikz-cd}
\newcommand{\Q}{\mathbb{Q}}
\newcommand{\C}{\mathbb{C}}
\lectureinfo{28}{Nov. 23, 2026}{Quotient Rings, the First Isomorphism Theorem, and the Correspondence Theorem}

% Typeset from handwritten Lecture 30, Nov. 24, 2025, pp. 1--11.
\begin{document}
\raggedbottom
\makelecturetitle

\section{Ideals and quotient rings}

Recall that an \emph{ideal} of a ring $R$ is a subset
$I\subseteq R$ such that
\begin{itemize}
  \item $0\in I$;
  \item $I$ is closed under addition;
  \item $ab\in I$ whenever $a\in I$ and $b\in R$.
\end{itemize}
From an ideal $I$ of $R$ we build the \emph{quotient ring} $R/I$.
Its elements are the cosets $a+I$ with $a\in R$, and the
multiplication is defined by
\[
  (a+I)(b+I)\coloneqq ab+I.
\]
The map
\[
  \pi:R\to R/I,\qquad a\mapsto a+I,
\]
is a ring homomorphism, often called the \emph{canonical quotient
map}.

\section{The universal property of quotients}

\begin{theorem}[Universal property of quotients]\label{thm:universal}
Let $\varphi:R\to S$ be a ring homomorphism and let $I$ be an ideal
of $R$ with $\varphi(I)=0$ in $S$. Then there is a unique ring
homomorphism $\tilde\varphi:R/I\to S$ making the following diagram
commute:
\[
  \begin{tikzcd}
    R \arrow[r,"\varphi"] \arrow[d,"\pi"'] & S\\
    R/I \arrow[ur,dashed,"\tilde\varphi"'] &
  \end{tikzcd}
\]
That is, $\tilde\varphi\circ\pi=\varphi$.
\end{theorem}
\begin{proof}
Commutativity forces the definition
\[
  \tilde\varphi(a+I)=\varphi(a),
\]
which gives uniqueness. This is well defined: if $a+I=a'+I$, then
$a'-a\in I$, so $\varphi(a')-\varphi(a)=\varphi(a'-a)=0$.
Hence $\varphi(a)=\varphi(a')$. Since $\varphi$ is a ring
homomorphism and the operations on $R/I$ are defined using
representatives, $\tilde\varphi$ is a ring homomorphism.
\end{proof}

% NOTE: p. 2 also has "Z[x]/(x^2+1)" written in the top margin with
% no further comment; omitted.

\section{The first isomorphism theorem for rings}

As a consequence, we obtain the first isomorphism theorem for rings.

\begin{remark}
A ring homomorphism $\varphi:R\to S$ is an \emph{isomorphism} if it
has an inverse map. Equivalently:
\begin{itemize}
  \item $\varphi$ is a bijective homomorphism; or
  \item $\varphi$ is a homomorphism with $\ker(\varphi)=0$ that is
        surjective; or
  \item $\varphi$ is a ring homomorphism that is an isomorphism of
        the underlying additive groups.
\end{itemize}
\end{remark}

\begin{theorem}[First isomorphism theorem for rings]\label{thm:first-iso}
If $\varphi:R\to S$ is a surjective homomorphism of rings, then
there is an isomorphism of rings
\[
  \bar\varphi:R/\ker(\varphi)\xrightarrow{\ \sim\ }S.
\]
\end{theorem}
\begin{proof}
Since $\varphi(\ker\varphi)=0$, Theorem~\ref{thm:universal} gives a
ring homomorphism $\bar\varphi$ fitting into the diagram
\[
  \begin{tikzcd}
    R \arrow[r,"\varphi"] \arrow[d] & S\\
    R/\ker(\varphi) \arrow[ur,dashed,"\bar\varphi"'] &
  \end{tikzcd}
\]
By the first isomorphism theorem for groups, applied to the
additive groups, $\bar\varphi$ is a bijection. A bijective ring
homomorphism is an isomorphism of rings.
\end{proof}

\section{The correspondence theorem}

As in group theory, we also have a correspondence theorem in ring
theory. Recall the group-theoretic version: if $\varphi:G\to H$ is a
surjective homomorphism with kernel $K=\ker(\varphi)$, there is a
bijective correspondence
\[
  \{\text{subgroups }L\text{ of }G\text{ with }L\supseteq K\}
  \longleftrightarrow
  \{\text{subgroups }L'\text{ of }H\}.
\]

\begin{theorem}[Correspondence theorem for rings]\label{thm:correspondence}
Let $\varphi:R\to S$ be a surjective homomorphism of rings and let
$K=\ker(\varphi)$.
\begin{enumerate}
  \item There is a bijective correspondence
  \[
    \{\text{ideals }I\text{ of }R\text{ with }I\supseteq K\}
    \longleftrightarrow
    \{\text{ideals }J\text{ of }S\},
  \]
  given by $I\mapsto\varphi(I)$, with inverse $J\mapsto\varphi^{-1}(J)$.
  \item If $I\leftrightarrow J$ under this correspondence, then
  \[
    R/I\cong S/J
  \]
  as rings.
\end{enumerate}
\end{theorem}
\begin{proof}
\emph{Preimages of ideals are ideals.}
Let $\varphi:R\to S$ be any ring homomorphism and let $J$ be an
ideal of $S$. We check that $\varphi^{-1}(J)$ is an ideal of $R$.
\begin{itemize}
  \item If $a,b\in\varphi^{-1}(J)$, then $\varphi(a),\varphi(b)\in J$,
  so $\varphi(a+b)=\varphi(a)+\varphi(b)\in J$. Thus
  $a+b\in\varphi^{-1}(J)$.
  \item $\varphi(0)=0\in J$, so $0\in\varphi^{-1}(J)$.
  \item If $a\in\varphi^{-1}(J)$ and $r\in R$, then
  $\varphi(ra)=\varphi(r)\varphi(a)\in J$, because $\varphi(r)\in S$,
  $\varphi(a)\in J$, and $J$ is an ideal of $S$. Thus
  $ra\in\varphi^{-1}(J)$.
\end{itemize}
Moreover, $\varphi^{-1}(J)\supseteq\varphi^{-1}(0)=K$.

\emph{Images of ideals are ideals.}
Now assume that $\varphi$ is surjective: for every $s\in S$ there is
$r\in R$ with $\varphi(r)=s$. Let $I$ be an ideal of $R$ with
$I\supseteq K$. We check that $\varphi(I)$ is an ideal of $S$.
It is closed under addition and contains $0$. For the absorption
property, let $a\in\varphi(I)$ and $s\in S$. Then $a=\varphi(b)$ for
some $b\in I$, and by surjectivity $s=\varphi(r)$ for some $r\in R$.
Hence
\[
  s\cdot a=\varphi(r)\varphi(b)=\varphi(r b)\in\varphi(I),
\]
since $rb\in I$. The rest of the verification that $\varphi(I)$ is
an ideal of $S$ is left as an exercise.

\emph{Bijectivity.}
Ideals $I\subseteq R$ and $J\subseteq S$ are in particular additive
subgroups, and $\varphi$ is a homomorphism of additive groups.
Since the correspondence theorem holds for groups, the maps
$I\mapsto\varphi(I)$ and $J\mapsto\varphi^{-1}(J)$ are mutually
inverse bijections; by the two steps above, they carry ideals to
ideals.

\emph{The quotients.}
Suppose $I\leftrightarrow J$, so $J=\varphi(I)$ and
$I=\varphi^{-1}(J)$. Consider the composite
\[
  \begin{tikzcd}
    R \arrow[r,"\varphi"] \arrow[d] & S \arrow[d,"\pi"]\\
    R/I \arrow[r,dashed] & S/J
  \end{tikzcd}
\]
The composite $\pi\circ\varphi$ is surjective, and it sends $I$ to
$0$ because $\varphi(I)=J$. By the universal property, it induces a
homomorphism $R/I\to S/J$. Its kernel consists of all $a$ with
$\varphi(a)\in J$, namely $\ker(\pi\circ\varphi)=\varphi^{-1}(J)=I$.
By the first isomorphism theorem, $R/I\cong S/J$.
\end{proof}

\begin{exercise}
Complete the proof that $\varphi(I)$ is an ideal of $S$ when
$\varphi:R\to S$ is surjective and $I$ is an ideal of $R$.
\end{exercise}

\section{Examples}

\begin{example}[Ideals of $\C{[t]}/(t^2-1)$]
Consider the canonical quotient map
\[
  \varphi:\C[t]\to\C[t]/(t^2-1).
\]
Every ideal of $\C[t]$ has the form $(f(t))$ for some polynomial $f$.
What are the ideals of $\C[t]/(t^2-1)$?

The map $\varphi$ is surjective with $\ker(\varphi)=(t^2-1)$.
By the correspondence theorem, ideals $J$ of $\C[t]/(t^2-1)$
correspond to ideals $I$ of $\C[t]$ containing $K=(t^2-1)$.
Writing $I=(f)$,
\[
  (f)\supseteq(t^2-1)
  \iff f\mid t^2-1=(t-1)(t+1)\text{ in }\C[t].
\]
Up to a nonzero scalar, $f$ is one of
\[
  1,\qquad t-1,\qquad t+1,\qquad (t-1)(t+1)=t^2-1.
\]
The correspondence is therefore
\[
\begin{array}{ccl}
  \text{Ideal of }\C[t] & & \text{Ideal of }\C[t]/(t^2-1)\\ \hline
  (1) & \longmapsto & (1)=\C[t]/(t^2-1)\\
  (t-1) & \longmapsto & \bigl(t-1+(t^2-1)\bigr)\\
  (t+1) & \longmapsto & \bigl(t+1+(t^2-1)\bigr)\\
  (t^2-1) & \longmapsto & 0
\end{array}
\]
% NOTE: the notes leave the images of (t-1) and (t+1) as "?"; the
% entries above are the images of the generators under phi.
\end{example}

\Needspace{16\baselineskip}
\begin{example}[Ideals of $\Q{[x]}/(x^3-2)$]
Consider the canonical quotient map
\[
  \varphi:\Q[x]\to\Q[x]/(x^3-2).
\]
What are the ideals of $\Q[x]/(x^3-2)$? They correspond to ideals
$I=(f)$ of $\Q[x]$ with
\[
  (f)\supseteq(x^3-2)\iff f\mid x^3-2\text{ in }\Q[x].
\]
Since $x^3-2$ does not factor over $\Q$, up to a nonzero scalar
either $f=1$ or $f=x^3-2$. The correspondence is
\[
  (1)\longmapsto(1),\qquad (x^3-2)\longmapsto(0).
\]
So $\Q[x]/(x^3-2)$ has only the two ideals $(0)$ and $(1)$, and it
is therefore a field! This field may be identified with
$\Q(\sqrt[3]{2})$.
% NOTE: p. 11 writes "Q(gamma)" without defining gamma; presumably
% gamma is a root of x^3-2, e.g. the real cube root of 2.
\end{example}

\end{document}
